% Generato dalla versione Word revisionata e dalla bozza auditata.
% Le citazioni numeriche originali sono state convertite in chiavi BibLaTeX.

\section{Dal vantaggio teorico al problema del rumore}\label{sec:finale_1_1}

La fattorizzazione di grandi numeri interi è un esempio di problema per il quale l'aumento della potenza di calcolo classica non modifica la natura della difficoltà computazionale: il miglior algoritmo classico generale noto, il General Number Field Sieve, ha complessità sub-esponenziale ma non polinomiale nella lunghezza dell'input \cite{buhler1993gnfs,rsa250_2020}. Nel 1994 Peter Shor ha mostrato che, nel modello ideale di calcolo quantistico, la fattorizzazione e il logaritmo discreto possono invece essere risolti in tempo polinomiale \cite{shor1994}. Il risultato ha una conseguenza immediata per la sicurezza informatica, perché sistemi a chiave pubblica come RSA fondano la propria sicurezza pratica proprio sulla difficoltà di fattorizzare interi di grandi dimensioni \cite{rsa1978}.

Il vantaggio algoritmico, tuttavia, non coincide con la possibilità di eseguire oggi Shor su istanze crittograficamente rilevanti. I processori quantistici disponibili operano in un regime nel quale errori di gate, decoerenza, readout imperfetto, connettività limitata e crosstalk riducono la profondità dei circuiti eseguibili con affidabilità \cite{preskill2018}. Shor è particolarmente sensibile a questi effetti: la ricerca del periodo richiede di preservare informazione di fase lungo l'esponenziazione modulare e di renderla osservabile attraverso la Quantum Phase Estimation (QPE) e la Quantum Fourier Transform (QFT). Se il rumore disperde la probabilità che dovrebbe concentrarsi intorno agli esiti legati al periodo, la successiva elaborazione classica può non ricostruire i fattori corretti. Studi recenti su order finding rumoroso mostrano inoltre che la recoverability dipende dalla struttura residua della distribuzione e non dalla sola presenza visiva dei picchi \cite{skosana2021,yang2026,smolin2013}. Dimostrazioni sperimentali su piccole istanze hanno inoltre mostrato che i blocchi essenziali di Shor possono essere realizzati su piattaforme fisiche molto diverse, dall\textquotesingle NMR agli ioni intrappolati \cite{vandersypen2001,monz2016}; tali prove restano però lontane da una fattorizzazione fault-tolerant su scala crittografica \cite{smolin2013,gidney2021}.

Questa distanza tra vantaggio teorico e realizzabilità costituisce il punto di partenza della tesi. Il problema non viene affrontato cercando una singola tecnica ``migliore'' in assoluto, ma studiando tre livelli diversi di intervento: usare meglio l'informazione rimasta nelle misure; proteggere l'informazione durante il calcolo; ridurre il numero di operazioni che espongono il circuito al rumore. A questi tre livelli si affianca un'analisi di sensibilità che misura come degrada una specifica implementazione di Shor al crescere di un errore per gate controllato. L'obiettivo è distinguere ciò che migliora realmente la metrica di interesse da ciò che aggiunge complessità senza un beneficio misurabile.

\section{Obiettivi e contributo del lavoro}\label{sec:finale_1_2}

Il lavoro adotta un principio metodologico comune a tutte le campagne: ogni tecnica viene confrontata con una baseline esplicita, sugli stessi dati quando il confronto lo consente, e con metriche coerenti con il problema osservato. Questa scelta è centrale perché, in presenza di rumore, un risultato apparentemente positivo può dipendere dal post-processing, dalla struttura particolare dell'istanza o da una metrica intermedia che non coincide con l'obiettivo finale. Tutti gli esperimenti della tesi sono simulativi; non vengono quindi presentati come una dimostrazione di fattibilità su hardware quantistico reale né come una previsione delle risorse necessarie per rompere chiavi crittografiche di dimensione reale.

La prima linea di ricerca riguarda il post-processing degli output rumorosi di Shor. Sull'istanza principale N = 15, a = 7, vengono confrontate tre strategie applicate agli stessi istogrammi: TOP-1, che verifica il solo esito più frequente; TOP-4, che verifica i quattro candidati più frequenti; e un metodo che usa un classificatore di machine learning per decidere quando attivare la stessa ricerca TOP-4. Il confronto tra il metodo appreso e TOP-4 senza filtro costituisce l'ablazione del classificatore. La SVM selezionata raggiunge elevate metriche predittive, ma non riduce il numero di iterazioni necessarie alla fattorizzazione rispetto alla baseline; TOP-4, invece, raggiunge mediamente il successo al primo tentativo nei due scenari considerati. Il risultato mostra che, in questo caso, il beneficio deriva dalla ricerca multi-candidato e non dal filtro appreso. Un confronto con la Zero-Noise Extrapolation (ZNE) completa la valutazione delle strategie a valle dell'esecuzione.

La seconda linea studia la Quantum Error Correction (QEC), cioè la protezione dell'informazione mentre il calcolo è in corso. Il percorso parte dal codice a ripetizione, prosegue con il codice di Steane \ensuremath{[[7,1,3]]} e arriva al surface code \cite{steane1996,fowler2012,googleQECBelowThreshold2025,googleSurfaceCode2023}. Nei circuiti di memoria viene misurato il tasso di errore logico al variare dell'errore fisico e, per il surface code, della distanza del codice. Le simulazioni riproducono la soppressione quadratica attesa per i codici a distanza tre e mostrano, nel surface code, il comportamento di soglia: sotto la soglia stimata aumentare la distanza riduce l'errore logico, mentre sopra soglia lo amplifica. Questi valori appartengono ai circuiti di memoria, al modello di rumore e ai decoder adottati e non vengono trasferiti direttamente alle porte logiche di Shor.

All'interno della QEC viene inoltre studiato un secondo impiego del machine learning. Invece di classificare l'istogramma finale, la rete riceve la sindrome di errore. Quando il modello di rumore corrisponde alle assunzioni del decoder analitico, una rete usata come sostituto non mostra un vantaggio stabile rispetto al Minimum Weight Perfect Matching (MWPM). In presenza di correlazioni non ben rappresentate dal grafo del matching, una strategia ibrida può invece migliorare la decisione analitica correggendone alcuni errori residui. Il confronto successivo con BP+OSD mostra però che una parte del margine può dipendere dalla qualità del decoder di riferimento. Il risultato non è quindi una superiorità generale del machine learning, ma individua il regime nel quale un modello appreso può essere utile: quando osserva struttura informativa che il metodo analitico non rappresenta adeguatamente.

La terza linea interviene direttamente sulla complessità del circuito mediante una QFT approssimata (AQFT), che elimina rotazioni controllate di piccolo angolo. L'approssimazione riduce porte e profondità, ma introduce una perdita di precisione: il problema consiste quindi nel verificare se il rumore evitato compensi l'errore algoritmico aggiunto. Nel pilota su Shor con N = 15, il grado di troncamento scelto sui dati di selezione e valutato su dati indipendenti aumenta il successo da 0,4792 a 0,6279, mentre le porte ECR a due qubit scendono da 362 a 279. Il miglioramento è favorito dall'ordine r = 4, che produce fasi esattamente rappresentabili con due bit, e non viene generalizzato automaticamente a istanze differenti. Per N = 21 il troncamento dell'aritmetica riduce invece le porte a due qubit da 49 661 a 23 178, ma diminuisce il successo ideale e nei test rumorosi a budget ridotto non emerge un beneficio statisticamente dimostrato. Nei benchmark separati di QPE, al contrario, la variante troncata risulta migliore in tutti i nove confronti effettuati, con sei intervalli di confidenza individuali interamente positivi. Questi risultati mostrano che l'efficacia dell'AQFT dipende dal compromesso tra precisione di fase, costo circuitale e regime di rumore \cite{barencoAQFT1996,akoramurthy2026}.

Un'analisi distinta misura infine la sensibilità di Shor a un errore Pauli uniforme applicato dopo le porte compilate. Il parametro di questo esperimento è un proxy per gate e non coincide con il tasso di errore logico per round o per memoria misurato nelle campagne QEC. Mantenere separate queste grandezze evita una conclusione non supportata: la tesi non simula un'esecuzione end-to-end di Shor fault-tolerant, ma caratterizza separatamente la robustezza dell'algoritmo, la capacità dei codici di sopprimere gli errori e il comportamento dei decoder.

Nel complesso, il contributo del lavoro è quindi soprattutto metodologico e sperimentale: costruire confronti appaiati e riproducibili, isolare mediante ablazione il contributo dei componenti, distinguere metriche che descrivono fenomeni diversi e verificare su dati indipendenti quando una maggiore complessità produce un beneficio reale. Il criterio che emerge è semplice: una tecnica aggiuntiva è utile quando sfrutta informazione che la baseline trascura oppure quando elimina un costo maggiore della perdita che introduce.

\section{Struttura della tesi}\label{sec:finale_1_3}

La tesi è organizzata in sei capitoli. Il Capitolo 2 formalizza il problema di ricerca, gli obiettivi, le domande sperimentali, le ipotesi e i criteri di valutazione. Il Capitolo 3 presenta i fondamenti di calcolo quantistico necessari alla lettura del lavoro, descrive l'algoritmo di Shor, i principali modelli di rumore e la metodologia sperimentale, limitando la teoria a ciò che serve per interpretare gli esperimenti. Il Capitolo 4 descrive l'implementazione, gli strumenti software, la costruzione e validazione dei circuiti, la generazione dei dati e il contratto di riproducibilità.

Il Capitolo 5 raccoglie i risultati: confronto TOP-1/TOP-4/machine learning e ZNE, codici a ripetizione e di Steane, surface code, decodifica analitica e appresa, analisi di sensibilità di Shor e studio della QFT approssimata. I risultati vengono discussi mantenendo separate le metriche dei diversi esperimenti e indicando per ciascuno il relativo perimetro di validità. Il Capitolo 6 sintetizza le evidenze, risponde alle domande di ricerca, esplicita i limiti del lavoro e individua gli sviluppi necessari per passare dalle simulazioni controllate a una validazione su hardware e, in prospettiva, a un'analisi fault-tolerant integrata.
